% --- Archon physics formalization source begin ---
% archon:physics
% archon:covers IPhO2026Problems/problem_IPhO_2026_3_A_1.lean
% archon:source-report reports/ipho_2026/problem_IPhO_2026_3_A_1.source.json
% archon:problem-id IPhO_2026_3
% archon:part-id A.1

\chapter{Physics problem IPhO\_2026\_3\_A\_1}
\label{ch:IPhO2026Problems_problem_IPhO_2026_3_A_1}

\paragraph{Problem source.}
A homogeneous isotropic paramagnetic torus has mean radius R, inner radius r
with r << R, volume V, and cross-sectional area A.  An insulated conducting
wire is wound densely around it with N turns and instantaneous current I.
Fields H and B and magnetization M are approximately uniform in the torus.
Use B = mu\_0*H + mu\_0*M, Ampere's law, and the sign convention that work and
heat entering the paramagnetic torus are positive.

Current subquestion:
Write the field magnitude H inside the torus in terms of N, I, A, and V.

\paragraph{Current subquestion.}
Write the field magnitude H inside the torus in terms of N, I, A, and V.

\paragraph{Recorded answer/context.}
H = N*I*A/V.

\paragraph{Figure/image path.}
/root/proposal\_for\_physic/science-mango/ipho\_2026\_source/image/T3\_page-2.png

\paragraph{Formalization target.}
create a compiling Lean file with sorry bodies at `IPhO2026Problems/problem\_IPhO\_2026\_3\_A\_1.lean`. The Lean declarations must preserve the physical quantities, dimensions or dimensional roles, figure labels, governing-law hypotheses, and final relation expressed by this problem.
Use Mathlib/Physlib names found through LeanExplore where available. If a physics API is missing, introduce faithful local abstractions rather than scalar placeholder aliases.

\begin{theorem}[Physics formalization target]
  \label{thm:physics:IPhO_2026_3_A_1:target}
  \lean{IPhO2026Problems.IPhO2026_3_A_1.fieldStrength_eq_turns_current_area_div_volume}
  \uses{def:physics:IPhO_2026_3_A_1:aux001, def:physics:IPhO_2026_3_A_1:aux002, def:physics:IPhO_2026_3_A_1:aux003, def:physics:IPhO_2026_3_A_1:aux004, def:physics:IPhO_2026_3_A_1:aux005, def:physics:IPhO_2026_3_A_1:aux006, def:physics:IPhO_2026_3_A_1:aux007, def:physics:IPhO_2026_3_A_1:aux008, def:physics:IPhO_2026_3_A_1:aux009, def:physics:IPhO_2026_3_A_1:aux010, def:physics:IPhO_2026_3_A_1:aux011, def:physics:IPhO_2026_3_A_1:aux012, def:physics:IPhO_2026_3_A_1:aux013, def:physics:IPhO_2026_3_A_1:aux014, def:physics:IPhO_2026_3_A_1:aux015, def:physics:IPhO_2026_3_A_1:aux016, def:physics:IPhO_2026_3_A_1:aux017, def:physics:IPhO_2026_3_A_1:aux018}
  For the Fig. 3a paramagnetic torus, Ampère's law and V = (2πR) A give H = N I A / V.
\end{theorem}
\begin{proof}
  Use the typed geometry, governing-law, branch, and measurement interfaces listed in the declaration index.  Eliminate the intermediate readouts in their physical order and then evaluate the requested exact or uncertainty relation.
\end{proof}
\section*{Formal declaration index}
The following blocks record the typed carriers and bridge declarations used by the target.  Their dependency lists follow the declaration-level model in the covered Lean file.

\begin{definition}[Declaration LengthMagnitude]
  \label{def:physics:IPhO_2026_3_A_1:aux001}
  \lean{IPhO2026Problems.IPhO2026_3_A_1.LengthMagnitude}
  A scalar length magnitude, retaining its physical dimension.
\end{definition}
\begin{proof}
  This is the typed definition of the stated physical carrier; its fields or defining equation provide the claimed data.
\end{proof}

\begin{definition}[Declaration AreaMagnitude]
  \label{def:physics:IPhO_2026_3_A_1:aux002}
  \lean{IPhO2026Problems.IPhO2026_3_A_1.AreaMagnitude}
  A scalar cross-sectional area magnitude, retaining its physical dimension.
\end{definition}
\begin{proof}
  This is the typed definition of the stated physical carrier; its fields or defining equation provide the claimed data.
\end{proof}

\begin{definition}[Declaration VolumeMagnitude]
  \label{def:physics:IPhO_2026_3_A_1:aux003}
  \lean{IPhO2026Problems.IPhO2026_3_A_1.VolumeMagnitude}
  A scalar volume magnitude, retaining its physical dimension.
\end{definition}
\begin{proof}
  This is the typed definition of the stated physical carrier; its fields or defining equation provide the claimed data.
\end{proof}

\begin{definition}[Declaration ElectricCurrentMagnitude]
  \label{def:physics:IPhO_2026_3_A_1:aux004}
  \lean{IPhO2026Problems.IPhO2026_3_A_1.ElectricCurrentMagnitude}
  The magnitude of an instantaneous electric current, with dimension charge/time.
\end{definition}
\begin{proof}
  This is the typed definition of the stated physical carrier; its fields or defining equation provide the claimed data.
\end{proof}

\begin{definition}[Declaration MagneticFieldStrengthMagnitude]
  \label{def:physics:IPhO_2026_3_A_1:aux005}
  \lean{IPhO2026Problems.IPhO2026_3_A_1.MagneticFieldStrengthMagnitude}
  The magnitude of H or M, with dimension current/length.
\end{definition}
\begin{proof}
  This is the typed definition of the stated physical carrier; its fields or defining equation provide the claimed data.
\end{proof}

\begin{definition}[Declaration MagneticFluxDensityMagnitude]
  \label{def:physics:IPhO_2026_3_A_1:aux006}
  \lean{IPhO2026Problems.IPhO2026_3_A_1.MagneticFluxDensityMagnitude}
  The magnitude of B, with its magnetic-flux-density dimension.
\end{definition}
\begin{proof}
  This is the typed definition of the stated physical carrier; its fields or defining equation provide the claimed data.
\end{proof}

\begin{definition}[Declaration MagneticPermeabilityMagnitude]
  \label{def:physics:IPhO_2026_3_A_1:aux007}
  \lean{IPhO2026Problems.IPhO2026_3_A_1.MagneticPermeabilityMagnitude}
  A magnetic permeability magnitude, including the vacuum permeability μ₀.
\end{definition}
\begin{proof}
  This is the typed definition of the stated physical carrier; its fields or defining equation provide the claimed data.
\end{proof}

\begin{definition}[Declaration EnergyMagnitude]
  \label{def:physics:IPhO_2026_3_A_1:aux008}
  \lean{IPhO2026Problems.IPhO2026_3_A_1.EnergyMagnitude}
  An energy magnitude, used only to record the problem's transfer sign convention.
\end{definition}
\begin{proof}
  This is the typed definition of the stated physical carrier; its fields or defining equation provide the claimed data.
\end{proof}

\begin{definition}[Declaration siValue]
  \label{def:physics:IPhO_2026_3_A_1:aux009}
  \lean{IPhO2026Problems.IPhO2026_3_A_1.siValue}
  The numerical readout of a dimensionful scalar quantity in SI units.
\end{definition}
\begin{proof}
  This is the typed definition of the stated physical carrier; its fields or defining equation provide the claimed data.
\end{proof}

\begin{definition}[Declaration apparatusFigureLabel]
  \label{def:physics:IPhO_2026_3_A_1:aux010}
  \lean{IPhO2026Problems.IPhO2026_3_A_1.apparatusFigureLabel}
  The official source labels the apparatus drawing as Fig. 3a.
\end{definition}
\begin{proof}
  This is the typed definition of the stated physical carrier; its fields or defining equation provide the claimed data.
\end{proof}

\begin{definition}[Declaration EnergyTransferDirection]
  \label{def:physics:IPhO_2026_3_A_1:aux011}
  \lean{IPhO2026Problems.IPhO2026_3_A_1.EnergyTransferDirection}
  Whether energy crosses the system boundary into or out of the Pm-T.
\end{definition}
\begin{proof}
  This is the typed definition of the stated physical carrier; its fields or defining equation provide the claimed data.
\end{proof}

\begin{definition}[Declaration signedEnergySI]
  \label{def:physics:IPhO_2026_3_A_1:aux012}
  \lean{IPhO2026Problems.IPhO2026_3_A_1.signedEnergySI}
  \uses{def:physics:IPhO_2026_3_A_1:aux008, def:physics:IPhO_2026_3_A_1:aux009, def:physics:IPhO_2026_3_A_1:aux011}
  The source convention: work and heat entering the paramagnetic torus are positive, while transfers leaving it are negative.
\end{definition}
\begin{proof}
  This is the typed definition of the stated physical carrier; its fields or defining equation provide the claimed data.
\end{proof}

\begin{definition}[Declaration HomogeneousIsotropicParamagneticTorus]
  \label{def:physics:IPhO_2026_3_A_1:aux013}
  \lean{IPhO2026Problems.IPhO2026_3_A_1.HomogeneousIsotropicParamagneticTorus}
  \uses{def:physics:IPhO_2026_3_A_1:aux001, def:physics:IPhO_2026_3_A_1:aux002, def:physics:IPhO_2026_3_A_1:aux003, def:physics:IPhO_2026_3_A_1:aux009}
  The homogeneous isotropic paramagnetic torus and its Fig. 3a geometry. minorRadius is the cross-sectional radius r marked by the diameter 2r in the figure. The explicit dimensionless bound records the approximation r ≪ R without choosing an unsupported numerical tolerance.
\end{definition}
\begin{proof}
  This is the typed definition of the stated physical carrier; its fields or defining equation provide the claimed data.
\end{proof}

\begin{definition}[Declaration DenseInsulatedWinding]
  \label{def:physics:IPhO_2026_3_A_1:aux014}
  \lean{IPhO2026Problems.IPhO2026_3_A_1.DenseInsulatedWinding}
  \uses{def:physics:IPhO_2026_3_A_1:aux004, def:physics:IPhO_2026_3_A_1:aux009}
  The densely wound insulated wire. Its instantaneous current is represented by its nonnegative magnitude, so no orientation branch is selected implicitly.
\end{definition}
\begin{proof}
  This is the typed definition of the stated physical carrier; its fields or defining equation provide the claimed data.
\end{proof}

\begin{definition}[Declaration UniformToroidalMagneticState]
  \label{def:physics:IPhO_2026_3_A_1:aux015}
  \lean{IPhO2026Problems.IPhO2026_3_A_1.UniformToroidalMagneticState}
  \uses{def:physics:IPhO_2026_3_A_1:aux005, def:physics:IPhO_2026_3_A_1:aux006, def:physics:IPhO_2026_3_A_1:aux009}
  The approximately uniform magnitudes of H, B, and M in the torus. For the isotropic paramagnet, the nonnegative scalar magnetization is the component along the same toroidal direction as fieldStrength.
\end{definition}
\begin{proof}
  This is the typed definition of the stated physical carrier; its fields or defining equation provide the claimed data.
\end{proof}

\begin{definition}[Declaration meanLoopLengthSI]
  \label{def:physics:IPhO_2026_3_A_1:aux016}
  \lean{IPhO2026Problems.IPhO2026_3_A_1.meanLoopLengthSI}
  \uses{def:physics:IPhO_2026_3_A_1:aux009, def:physics:IPhO_2026_3_A_1:aux013}
  The length 2πR of the circular Ampèrian loop through the torus.
\end{definition}
\begin{proof}
  This is the typed definition of the stated physical carrier; its fields or defining equation provide the claimed data.
\end{proof}

\begin{definition}[Declaration ParamagneticConstitutiveLaw]
  \label{def:physics:IPhO_2026_3_A_1:aux017}
  \lean{IPhO2026Problems.IPhO2026_3_A_1.ParamagneticConstitutiveLaw}
  \uses{def:physics:IPhO_2026_3_A_1:aux007, def:physics:IPhO_2026_3_A_1:aux009, def:physics:IPhO_2026_3_A_1:aux015}
  The scalar consequence of B = μ₀ H + μ₀ M for the common toroidal direction. It is recorded as a governing law, not as a definition of any field.
\end{definition}
\begin{proof}
  This is the typed definition of the stated physical carrier; its fields or defining equation provide the claimed data.
\end{proof}

\begin{definition}[Declaration ToroidalAmpereLaw]
  \label{def:physics:IPhO_2026_3_A_1:aux018}
  \lean{IPhO2026Problems.IPhO2026_3_A_1.ToroidalAmpereLaw}
  \uses{def:physics:IPhO_2026_3_A_1:aux009, def:physics:IPhO_2026_3_A_1:aux013, def:physics:IPhO_2026_3_A_1:aux014, def:physics:IPhO_2026_3_A_1:aux015, def:physics:IPhO_2026_3_A_1:aux016}
  Ampère's circuital law reduced using the approximately uniform toroidal field: the circulation H (2πR) equals the enclosed free current N I.
\end{definition}
\begin{proof}
  This is the typed definition of the stated physical carrier; its fields or defining equation provide the claimed data.
\end{proof}
% --- Archon physics formalization source end ---
